% META: {"id": "1SPE-VARALEA-EX-007", "chapitre": "1SPE-VARIABLES-ALEATOIRES", "type_objet": "exercice", "capacites": ["1SPE-VARIABLES-ALEATOIRES-C1"], "capacites_codes": ["C1"], "parcours": 2, "competences": ["calculer", "raisonner"], "duree_min": 15, "mode_creation": "ex_nihilo", "sources_inspiration": [], "fichier_tex": "chapitres/1SPE-VARIABLES-ALEATOIRES/exercices/1SPE-VARALEA-EX-007.tex", "corrige_tex": "chapitres/1SPE-VARIABLES-ALEATOIRES/corriges/1SPE-VARALEA-CO-007.tex", "status": "generated"}
% BEGIN-VERIFY
% from sympy import Rational
% # Tirage sans remise 2 boules parmi 3R+2B, X = nb rouges
% # Calcul par l'arbre ordonne des deux tirages sans remise
% p0 = Rational(2,5)*Rational(1,4)
% p1 = Rational(3,5)*Rational(2,4) + Rational(2,5)*Rational(3,4)
% p2 = Rational(3,5)*Rational(2,4)
% assert p0 == Rational(1,10)
% assert p1 == Rational(6,10)
% assert p2 == Rational(3,10)
% assert p0+p1+p2 == 1
% END-VERIFY
\begin{exercice}{1SPE-VARALEA-EX-007}{2}{15}
Une urne contient $3$ boules rouges et $2$ boules bleues. On tire $2$ boules au hasard \emph{sans remise}. Soit $X$ le nombre de boules rouges obtenues.

\begin{enumerate}
  \item Quelles sont les valeurs possibles de $X$ ?
  \item Construire l'arbre des deux tirages, puis calculer $P(X=0)$, $P(X=1)$ et $P(X=2)$.
  \item Dresser le tableau de loi de $X$ et vérifier.
\end{enumerate}
\end{exercice}
